平行板电容器有效重叠面积为 \(A\)，板距为 \(d>0\)，忽略边缘效应，填充均匀理想线性介质介电常数 \(\varepsilon>0\)，电容 \(C(d)=\varepsilon A/d\)。两板可在法向缓慢移动，静电场始终准静态。比较断开电源保持正电荷量 \(Q\) 不变以及接上理想恒压电源保持电压 \(V\) 不变这两种控制方式。
(1) 在固定 \(Q\) 下求电场能及板间吸引力，明确力方向与对间距的导数。
(2) 在恒压 \(V\) 下同时求场能变化、电源做功及机械做功，解释为何不能只微分场能求力。
(3) 从电源—电容器整体势能构造适于固定电压的有效自由能，独立推导极板受力。
(4) 检验两种条件下瞬时电压力可写成同一电场强度形式，并分析缓慢拉开极板时能量流向。
